\section{Voting While the Previous Epoch Closes}
\label{sec:pipeline}
A validator opens epoch $e+1$ the moment it leaves epoch $e$, typically a few seconds before $S_e$ is decided. During this overlap two things could go wrong: the new epoch could confirm a block that the old epoch's latent votes contradict, and the checkpoint could later rule against something the new epoch already confirmed. Rules~\ref{r:gate} to~\ref{r:cross} prevent both. We explain them with Alice's account at position 8, right after $A$ became final at position~7.

\para{The gate (\ref{r:gate}).}
The simplest safe behaviour is to wait: a block of epoch $e+1$ collects first and final votes as usual, but no validator treats it as final until it has installed $S_e$ and re-checked the block against it. If Alice's block $B$ at position 8 is published during the overlap, its certificate is held, and the moment $S_e$ arrives the held certificate takes effect, or is discarded if $S_e$ kept a conflicting block. This is correct but it is also the stop-and-drain cost in disguise: every block published during a close waits for the close to end. Our measurements of the gate alone showed confirmations stopping for one to three seconds at every boundary (supplement).

\para{The overlap rule (\ref{r:overlap}).}
A block can skip the gate if the old epoch provably cannot contradict it. The proof is an exclusion witness in the closing epoch: four supporters for $B$ in epoch $e$. If four validators supported $B$ at position 8 in epoch $e$, no rival at position 8 can be notarized or final in epoch $e$ (Section~\ref{sec:epoch}), so $S_e$ cannot finalize a rival, cannot lock a rival for notarization, and any recovery lock $S_e$ places on a rival has origin $e$ and is released by exactly this witness. So a validator that holds the epoch-$e$ witness for $B$, four epoch-$(e+1)$ first votes for $B$, and an epoch-$(e+1)$ final certificate for $B$ may finalize $B$ before $S_e$ is known. The same check must hold for every unresolved block below $B$ on Alice's chain, because finalizing $B$ also finalizes its ancestors; and if $B$ bypasses a recovery lock inherited from an older checkpoint, the validator must hold a witness of that lock's own origin epoch. The rule is exactly the check a validator would perform after installing $S_e$, so a validator that finalizes $B$ before and one that re-checks after reach the same verdict (Lemma~\ref{lem:verdict}).

\para{Late notarization (\ref{r:late}).}
A block published after the boundary has no epoch-$e$ votes at all, so the overlap rule would never apply to it and it would wait at the gate. Late notarization fixes this. A validator that has left epoch $e$ but not yet installed $S_e$, when it first-votes $B$ in epoch $e+1$, also signs an epoch-$e$ \emph{support} for $B$, unless it supported a different block at position 8 in epoch $e$. These late supports count toward the exclusion witness and toward nothing else: no notarization, no final vote, no report entry rests on them, since reports are frozen before they are cast. Safety is the same argument as before: one validator supports one block per position in epoch $e$, by a first vote or by a late support. In our runs 80 to 97\% of the blocks published during a close reached their witness this way and finalized before the checkpoint.

\para{Early votes (\ref{r:early}).}
Here is where the first of our earlier rules failed. Suppose $B$ conflicts with a block $B'$ that $S_e$ keeps under a recovery lock, and $V_3,V_4,V_5$ first-vote $B$ early, without the epoch-$e$ witness for $B$. $V_1$ later installs $S_e$, obtains the witness from the Byzantine $V_6$, releases the lock, and first-votes $B$ too: five first votes, a fast certificate that only $V_1$ can justify. The next checkpoint may select reports without $V_1$'s and keep $B'$, and a later epoch finalizes $B'$ while $B$'s certificate is still valid (supplement). The rule removes the dependency: a first vote names the checkpoint its signer had installed; \emph{early} votes, cast before $S_e$, count toward notarization and witnesses but never toward a fast certificate, and an early final vote is released only by a validator that holds the overlap witnesses, which its report then carries to the next checkpoint. A block first-voted during a close therefore takes two rounds instead of one.

\para{No switching (\ref{r:cross}).}
A validator that supported $B'$ at position 8 in epoch $e$ does not first-vote a conflicting $B$ in epoch $e+1$ until it has installed $S_e$. Otherwise it could contribute to an epoch-$(e+1)$ notarization of $B$ while its epoch-$e$ vote still helps $B'$ toward a latent certificate. At a position where it has seen two blocks, or which a checkpoint keeps, it casts no first vote at all before installing $S_e$: re-voting its old side early would let the split repeat in every epoch.

\para{After installation.}
When $S_e$ arrives, each validator re-checks its provisional epoch-$(e+1)$ work against the finalized history, the frontier, and the locks of $S_e$. Held certificates take effect; work on a block that $S_e$ excluded is discarded, without clearing the validator's recorded first vote. Every position with one kept block gets an instance at that position in the open epoch, and a rival is refused while any lock on the kept block stands. If a kept block does not finalize in epoch $e+1$ either, its first votes are exposed again in the next reports, and the next checkpoint locks it again with origin $e+1$: a block may carry locks of several origins, each released only by a witness of its own epoch. From then on every validator that installed the checkpoint votes it there.
